\section{Question 1: Methods and Implementation}

\subsection{Preprocessing and Bag-of-Words}

The first stage uses \texttt{emails\_1.csv} to demonstrate preprocessing and sparse lexical representation. Labels are normalized by mapping \texttt{spam} $\rightarrow 1$ and \texttt{ham} $\rightarrow 0$. Text is cleaned by removing non-alphabetic symbols and converting to lowercase. A CountVectorizer with \texttt{max\_features=15} is then used to construct a compact Bag-of-Words (BoW) representation.

This stage serves two roles: (i) verifying the preprocessing logic required by the assignment, and (ii) generating interpretable lexical statistics for report visualization.

\subsection{Training Data for Model Learning}

For model training, we use \texttt{emails\_2.csv}, which already contains BoW-style numeric features at scale. The pipeline performs the following transformations:
\begin{enumerate}
    \item Read CSV from fixed project path.
    \item Use \texttt{Prediction} (or \texttt{status} if present) as the binary target.
    \item Drop non-feature identifiers such as \texttt{Email No.}.
    \item Coerce all feature columns to numeric and fill missing values with 0.
\end{enumerate}

The final matrix has 5172 samples with a class distribution of 3672 ham and 1500 spam, preserving moderate imbalance representative of operational email streams.

\subsection{Train/Test Split and Reproducibility}

A fixed random seed is used with an 80/20 train-test split. Stratification is applied when class counts permit safe stratified partitioning. A helper guard disables stratification only in edge cases where a fallback mini-dataset would otherwise make stratified splitting invalid. This avoids silent runtime crashes while preserving expected behavior on the real dataset.

\subsection{Manual Evaluation Metrics}

The assignment requires manual implementations of binary classification metrics. Let TP, TN, FP, and FN denote confusion matrix counts. The metric definitions are:
\begin{align}
\text{Accuracy} &= \frac{TP + TN}{TP + TN + FP + FN} \\
\text{Precision} &= \frac{TP}{TP + FP} \\
\text{Recall} &= \frac{TP}{TP + FN} \\
F_1 &= \frac{2 \cdot \text{Precision} \cdot \text{Recall}}{\text{Precision} + \text{Recall}}
\end{align}

Zero-division safeguards are included for edge conditions.

\subsection{From-Scratch Logistic Regression}

Logistic Regression is implemented with gradient descent on binary cross-entropy loss:
\begin{align}
\hat{y} &= \sigma(w^Tx + b), \quad \sigma(z)=\frac{1}{1+e^{-z}} \\
\mathcal{L} &= -\frac{1}{N}\sum_{i=1}^{N}\left[y_i\log(\hat{y}_i)+(1-y_i)\log(1-\hat{y}_i)\right]
\end{align}

The implementation includes:
\begin{itemize}
    \item Feature normalization using training-set mean and standard deviation.
    \item Numerically stable sigmoid via clipping.
    \item Fixed learning rate and iteration count.
\end{itemize}

This model is discriminative and can capture correlated feature effects, which is often beneficial in sparse lexical spaces \cite{hastie2009elements,manning2008introduction}.

\subsection{From-Scratch Multinomial Naive Bayes}

Multinomial Naive Bayes is implemented with Laplace smoothing. For class $c$ and feature $j$:
\begin{align}
P(x_j|c) = \frac{N_{jc}+\alpha}{\sum_k N_{kc}+\alpha d}
\end{align}
where $d$ is feature dimension and $\alpha=1.0$. Training uses log probabilities to avoid underflow. Prediction computes:
\begin{align}
\log P(c|x) \propto \log P(c) + \sum_j x_j \log P(x_j|c)
\end{align}

Despite its conditional independence assumption, Multinomial NB remains a strong lexical baseline \cite{mccallum1998comparison}.

\subsection{Reporting Figures for Q1}

To support direct report integration, the notebook exports Q1 figures automatically:
\begin{itemize}
    \item Confusion matrix comparison.
    \item ROC and Precision-Recall curves for Logistic Regression.
    \item Top BoW words per class from \texttt{emails\_1.csv}.
    \item Metric bar chart comparing Logistic Regression vs Naive Bayes.
\end{itemize}

These files are saved under \texttt{answer/report/figures/} with deterministic names.
